% sec-04-trial-protocol.tex - Section 4, Trial Protocol.  A main section.
% FINAL stage.  Figures 10, 11, 12 and 13 drawn; Tables 11, 12 and 13 populated
% from trial-protocol/final-protocol/publication and
% trial-phase-2/final-protocol/publication/author.  Senior-author pass: Table 13
% loses its Characters column padding to the Elapsed column, and two closing
% subsections state the safety architecture in one paragraph and name the
% instruments the two studies would need to open.
\section{Trial Protocol}
\label{sec:protocol}

\subsection{The Phase 1 study}

The Phase 1 protocol, version 1.0.0, deposited June 21, 2026, describes itself
as an ``open-label, single-arm, dose-finding study using a standard 3+3
escalation across three daraxonrasib dose levels (DL1 through DL3) with
staggered (sentinel) enrollment'', planned for up to n = 18 participants
\cite{onpremwhippl}. A mandatory gate precedes any patient-facing robotic
activity: ``a mandatory Phase 0 simulation-validation gate ($\geq 1000$
simulated procedures across $\geq 2$ independent frameworks, with a Unified
Safety Level (USL) rating $\geq 7.0$) precedes any patient-facing robotic
activity, and prespecified pause and stopping rules govern progression''
\cite{onpremwhippl}. The intervention is stated with the same numeric
precision: perioperative daraxonrasib with a planned preoperative pause and an
LLM-advised postoperative restart at T+7, T+14, or T+21 days keyed to fistula
status and drug trough, combined with an eight-arm robotic Whipple performed
under continuous human oversight as a Class II collaborative device, with a 10
kHz heartbeat broadcast bus, a cross-arm emergency-stop budget at or below 3 ms,
per-arm tip-force caps at or below 3 N, a cumulative cross-arm cap at or below
18 N, and vascular no-fly gating around the superior mesenteric and portal veins
\cite{onpremwhippl}. Table 11 is the deposited synopsis in full.

\begin{apptable}
\begin{tabularx}{\textwidth}{>{\raggedright\arraybackslash}p{3.15cm} >{\raggedright\arraybackslash}X}
\headrow \thead{Element} & \thead{Phase 1, first-in-human, v1.0.0, June 21, 2026}\\
\midrule
Design & Open-label, single-arm, 3+3 escalation, DL1 to DL3, staggered sentinel enrollment \\
Sample size & Up to n = 18 \\
Primary endpoints & Dose-limiting toxicity rate during escalation; device- and procedure-related serious adverse event rate through day 30, including Clavien-Dindo grade III or higher \\
Key secondary & R0 rate at day-90 pathology; ISGPS grade B or C fistula rate; 90-day mortality; progression-free survival; overall survival to 24 months \\
Population & Adults 18 or older, resectable or borderline-resectable KRAS G12-mutated PDAC, ECOG 0 or 1, adequate organ function \\
Gate & Phase 0: at least 1000 simulations across at least 2 frameworks, Unified Safety Level at or above 7.0 \\
Device limits & 10 kHz heartbeat bus; cross-arm stop at or below 3 ms; per-arm tip force at or below 3 N; cross-arm at or below 18 N \\
Duration & Up to 28 days screening, surgery on day 0, day 1 to 7 acute window, day 30 primary safety window, day 90 pathology, follow-up every 12 weeks to 24 months \\
Regulatory basis & IND 21 CFR \S 312 plus IDE 21 CFR \S 812, Physical AI Subpart J \\
\end{tabularx}
\end{apptable}
\vspace{-0.6cm}
\tabcap{Table 11. The Phase 1 protocol synopsis as deposited, with the gate that\\must
close before any patient-facing robotic activity begins.}

Figure 10 takes the participant's point of view rather than the study's. Nine
forward states run down a left column at a constant 1.30 cm pitch, two exception
states sit in a right column 8.20 cm away, and each of the fourteen transitions
carries a guard written as a quantity that can be evaluated at the bedside. The
guards leaving any one state partition the space: the complement of ``KRAS G12
present and ECOG 0 or 1'' is the screen-failure guard, the complement of
``resection complete, no intraoperative stop'' is the force-cap guard, and the
complement of ``fistula grade below B and drug trough in range'' is the
dose-limiting toxicity guard. No participant can therefore occupy two successors
at once, and no state can deadlock.

% ---------------------------------------------------------------------------
% FIGURE 10.  plantuml-type, state with guards.  Nine forward states at 1.30 cm
% in a left column, two exception states in a right column, one dashed column
% rule at x = 7.30, and one upward return edge routed through a waypoint.
% ---------------------------------------------------------------------------
\begin{appfloat}[!tb]
\begin{appfig}[plantuml-type / state with guards]
\draw[trialslate,dashed,line width=0.4pt] (7.30,0.95) -- (7.30,-12.10);
\node[font=\tiny\sffamily,text=trialslate,rotate=90,anchor=south] at (7.20,-6.00) {exception column};
\node[umlinit] (i0) at (3.20,0.55) {};
\node[umlstatesoft,text width=30mm] (st1) at (3.20,-0.30) {Screening};
\node[umlstate,text width=30mm] (st2) at (3.20,-1.60) {Consented};
\node[umlstate,text width=30mm] (st3) at (3.20,-2.90) {Dose assigned};
\node[umlstateon,text width=30mm] (st4) at (3.20,-4.20) {Operative};
\node[umlstate,text width=30mm] (st5) at (3.20,-5.50) {Acute window};
\node[umlstatesoft,text width=30mm] (st6) at (3.20,-6.80) {Restart advised};
\node[umlstate,text width=30mm] (st7) at (3.20,-8.10) {Day 30 safety};
\node[umlstate,text width=30mm] (st8) at (3.20,-9.40) {Day 90 pathology};
\node[umlstateon,text width=30mm] (st9) at (3.20,-10.70) {Long term follow up};
\node[umlfinal] (fin) at (3.20,-11.70) {};
\fill[trialburg] (fin.center) circle (1.1mm);
\node[umlstategray,text width=30mm] (sf) at (11.40,-0.30) {Screen failure};
\node[umlstategray,text width=30mm,line width=0.9pt] (hd) at (11.40,-4.85) {Hold};
\node[umlfinal] (hf) at (11.40,-6.60) {};
\fill[trialburg] (hf.center) circle (1.1mm);
\draw[umlarrow] (i0) -- (st1);
\draw[umlarrow] (st1) -- node[umlguard,pos=0.5] {[KRAS G12 and ECOG 0 or 1]} (st2);
\draw[umlarrow] (st2) -- node[umlguard,pos=0.5] {[Phase 0 closed, USL 7.0 or above]} (st3);
\draw[umlarrow] (st3) -- node[umlguard,pos=0.5] {[cohort slot open, sentinel elapsed]} (st4);
\draw[umlarrow] (st4) -- node[umlguard,pos=0.5] {[resection complete, no stop]} (st5);
\draw[umlarrow] (st5) -- node[umlguard,pos=0.5] {[fistula below B, trough in range]} (st6);
\draw[umlarrow] (st6) -- node[umlguard,pos=0.5] {[restart at T+7, T+14 or T+21]} (st7);
\draw[umlarrow] (st7) -- node[umlguard,pos=0.5] {[no new grade III or higher]} (st8);
\draw[umlarrow] (st8) -- node[umlguard,pos=0.5] {[pathology adjudicated, R0 recorded]} (st9);
\draw[umlarrow] (st9) -- node[umlguard,pos=0.5] {[24-month OS reached]} (fin);
\draw[umlarrow] (st1.east) -- node[umlguard,pos=0.5] {[KRAS G12 absent or ECOG above 1]} (sf.west);
\draw[umlarrow] (st4.east) -- node[umlguard,pos=0.45] {[tip force above 3 N or cross arm above 18 N]} (hd.north west);
\draw[umlarrow] (st5.east) -- node[umlguard,pos=0.45] {[DLT within the day 30 window]} (hd.south west);
\draw[umlarrow] (hd.north) -- ++(0,-0.0) -- (11.40,-2.90) -- node[umlguard,pos=0.55] {[de escalation approved]} (st3.east);
\draw[umlarrow] (hd.south) -- node[umlguard,pos=0.5] {[stopping rule met]} (hf);
\pnote{0}{-12.60}{Nine forward states at one 1.30 cm pitch read as a clock; the two exception
states sit 8.20 cm away across a\\dashed column rule. The return edge rises to
y = -2.90 through a waypoint 1.30 cm clear of the Operative state.}
\end{appfig}
\vspace{-0.6cm}
\figcaption{Figure 10. One participant from screening to the 24-month endpoint,
with the\\guard on every transition and the two conditions that route to a study
hold.}
\end{appfloat}

\subsection{The Phase 2 study}

The Phase 2 protocol, version 1.1.0, deposited June 23, 2026, two days after its
predicate, is a ``multicenter, randomized 1:1, parallel-group, controlled,
open-label study with BICR of imaging, central masked pathology, and blinded
endpoint adjudication, conducted at eight high-volume HPB centers with a planned
sample size of n = 220 randomized participants'' \cite{phase2proto}. It states
its own dependence on the predicate plainly: the first-in-human study
``established the daraxonrasib RP2D and demonstrated the feasibility and hard
cyber-physical safety of an on-premises LLM-directed robotic Whipple, but it was
not powered for efficacy and explicitly deferred the comparative question to a
randomized trial'' \cite{phase2proto}. The trial targets approximately 140
progression-free survival events for 85 percent power at a two-sided alpha of
0.05 to detect a hazard ratio of 0.60, median PFS 14.0 months control against
23.3 months experimental, with a single interim analysis at approximately 60
percent of events using O'Brien-Fleming spending. Table 12 is the Phase 2
synopsis, and every randomized element in it depends on a Phase 1 result.

\begin{apptable}
\begin{tabularx}{\textwidth}{>{\raggedright\arraybackslash}p{3.15cm} >{\raggedright\arraybackslash}X}
\headrow \thead{Element} & \thead{Phase 2, randomized, v1.1.0, June 23, 2026}\\
\midrule
Design & Multicenter, randomized 1:1, parallel-group, controlled, open-label, with blinded independent central review \\
Sample size & n = 220, approximately 110 per arm, eight high-volume HPB centers \\
Primary endpoint & Progression-free survival per RECIST 1.1, investigator assessed, BICR as sensitivity analysis \\
Key secondary & Overall survival; R0 rate; ISGPS grade B or C fistula rate; major pathologic response; KRAS ctDNA clearance at week 12 \\
Statistics & About 140 PFS events, 85 percent power, two-sided alpha 0.05, target hazard ratio 0.60, median 14.0 against 23.3 months \\
Interim analysis & One, at about 60 percent of events, O'Brien-Fleming spending, independent monitoring board \\
Stratification & Resectability, KRAS allele G12D against other G12, neoadjuvant therapy, site \\
Predicate & Phase 1 v1.0.0, which established the RP2D and the safety architecture \\
\end{tabularx}
\end{apptable}
\vspace{-0.6cm}
\tabcap{Table 12. The Phase 2 protocol synopsis as deposited, and the Phase 1\\predicate
whose result each randomized element depends on.}

Figure 11 draws the clinical progression that joins the two documents, and
Figure 12 draws the documentary one. Four of the twelve Phase 1 sections were
carried into Phase 2 unchanged, four were replaced outright because the study is
no longer single-arm, and four elements exist in the Phase 2 document only
because randomization requires them. Table 13 states the elapsed production
time for both.

% ---------------------------------------------------------------------------
% FIGURE 11.  mermaid-type, flowchart LR.  Five columns, Phase 1 rungs at a
% 1.40 cm pitch inside a tighter band than the Phase 0 and Phase 2 pairs.
% ---------------------------------------------------------------------------
\begin{appfloat}[!tb]
\begin{appfig}[mermaid-type / flowchart LR]
\node[mmsoft,text width=28mm] (s1) at (0,1.05) {Simulation validation\\at least 1000 procedures\\across 2 frameworks};
\node[mmsoft,text width=28mm] (s2) at (0,-0.75) {Unified Safety Level\\rating at or above 7.0};
\begin{scope}[on background layer]\node[mmlane,fit=(s1)(s2)] (p0) {};\end{scope}
\node[mmlanetitle,anchor=south west] at ([yshift=1.3mm]p0.north west) {Phase 0, before any patient contact};
\node[mmdec,aspect=1.9,text width=20mm] (g0) at (4.05,0.15) {Gate 0 both hold?};
\node[mmmid,text width=28mm] (d1) at (7.90,1.55) {DL1 cohort\\3+3, sentinel staggered};
\node[mmmid,text width=28mm] (d2) at (7.90,0.15) {DL2 cohort\\3+3};
\node[mmmid,text width=28mm] (d3) at (7.90,-1.25) {DL3 cohort\\3+3, RP2D declared};
\begin{scope}[on background layer]\node[mmlane,fit=(d1)(d2)(d3)] (p1) {};\end{scope}
\node[mmlanetitle,anchor=south west] at ([yshift=1.3mm]p1.north west) {Phase 1, first in human, n up to 18};
\node[mmdec,aspect=1.9,text width=20mm] (g1) at (11.60,0.15) {Gate 1 DLT below limit?};
\node[mmgray,text width=27mm] (hold) at (11.60,-2.85) {Stop or de escalate\\prespecified pause rule};
\node[mmsoft,text width=28mm] (r1) at (15.20,1.15) {Randomize 1 to 1\\8 high volume centers};
\node[mmsoft,text width=28mm] (r2) at (15.20,-0.35) {PFS primary\\about 140 events};
\begin{scope}[on background layer]\node[mmlane,fit=(r1)(r2)] (p2) {};\end{scope}
\node[mmlanetitle,anchor=south west] at ([yshift=1.3mm]p2.north west) {Phase 2, multicenter randomized, n = 220};
\node[mmgoal,text width=32mm] (out) at (15.20,-2.35) {Hazard ratio 0.60 target\\85 percent power\\two sided alpha 0.05};
\draw[mmedge] (s1.east) -- (g0.north west);
\draw[mmedge] (s2.east) -- (g0.south west);
\draw[mmedgeb] (g0.east) -- node[mmlabel,above] {yes} (d2.west);
\draw[mmedge] (d1) -- (d2);
\draw[mmedge] (d2) -- (d3);
\draw[mmedge] (d3.east) -- (g1.south west);
\draw[mmedgeb] (g1.east) -- node[mmlabel,above] {yes} (r2.west);
\draw[mmedged] (g1.south) -- node[mmlabel,right] {no} (hold.north);
\draw[mmedged] (g0.south) -- ++(0,-1.35) -- (4.05,-3.55) -- (11.60,-3.55) -- node[mmlabel,pos=0.4] {no} (hold.south);
\draw[mmedge] (r1) -- (r2);
\draw[mmedgeb] (r2) -- (out);
\pnote{0}{-4.35}{The Phase 1 rungs sit at a 1.40 cm pitch and the Phase 0 and Phase 2 pairs at
1.80 and 1.50 cm, so the escalation\\band reads denser than the bands either side
of it. The two hold edges enter at different anchors.}
\end{appfig}
\vspace{-0.6cm}
\figcaption{Figure 11. From the Phase 0 simulation gate through three 3+3 dose
levels to\\Phase 2 randomization, with the quantity that must hold on each rung
shown.}
\end{appfloat}

% ---------------------------------------------------------------------------
% FIGURE 12.  d2-type, containers.  Nesting depth two.  Three equal-width lower
% containers at a 0.50 cm gutter against a 0.20 cm inter-leaf gutter.
% ---------------------------------------------------------------------------
\begin{appfloat}[!tb]
\begin{appfig}[d2-type / containers]
\foreach \i/\lab in {0/{Compliance},1/{Summary and schema},2/{Introduction},
  3/{Objectives},4/{Design},5/{Population}}{
  \node[d2gray,text width=20mm,minimum height=0.60cm,anchor=north west]
    at ({0.30+2.38*\i},-0.55) {\lab};}
\foreach \i/\lab in {0/{Intervention},1/{Discontinuation},2/{Assessments},
  3/{Statistics},4/{Oversight},5/{Additional}}{
  \node[d2gray,text width=20mm,minimum height=0.60cm,anchor=north west]
    at ({0.30+2.38*\i},-1.35) {\lab};}
\begin{scope}[on background layer]
\node[d2cont,minimum width=14.6cm,minimum height=2.05cm,anchor=north west] (pc) at (0,0) {};
\end{scope}
\node[d2title,anchor=north west] at ([shift={(2mm,-1.5mm)}]pc.north west) {Phase 1 protocol, v1.0.0, June 21 2026};
% carried
\foreach \i/\lab in {0/{Compliance},1/{Introduction},2/{Discontinuation},3/{Additional}}{
  \node[d2soft,text width=33mm,minimum height=0.62cm,anchor=north west]
    at (0.30,{-3.85-0.70*\i}) {\lab};}
\begin{scope}[on background layer]
\node[d2cont,minimum width=4.55cm,minimum height=3.15cm,anchor=north west] (cc) at (0,-3.20) {};\end{scope}
\node[d2title,anchor=north west] at ([shift={(2mm,-1.5mm)}]cc.north west) {Carried unchanged};
% replaced
\foreach \i/\lab in {0/{Summary and schema},1/{Objectives and endpoints},2/{Design},3/{Statistics}}{
  \node[d2gray2,text width=33mm,minimum height=0.62cm,anchor=north west]
    at (5.35,{-3.85-0.70*\i}) {\lab};}
\begin{scope}[on background layer]
\node[d2cont2,minimum width=4.55cm,minimum height=3.15cm,anchor=north west] (rc) at (5.05,-3.20) {};\end{scope}
\node[d2title2,anchor=north west] at ([shift={(2mm,-1.5mm)}]rc.north west) {Replaced outright};
% added
\foreach \i/\lab in {0/{Blinded independent central review},1/{Stratified permuted block randomization},
  2/{Interim analysis and alpha spending},3/{Multicenter site qualification}}{
  \node[d2key,text width=33mm,minimum height=0.62cm,anchor=north west]
    at (10.40,{-3.85-0.70*\i}) {\lab};}
\begin{scope}[on background layer]
\node[d2cont,minimum width=4.50cm,minimum height=3.15cm,anchor=north west,line width=0.9pt] (ac) at (10.10,-3.20) {};\end{scope}
\node[d2title,anchor=north west] at ([shift={(2mm,-1.5mm)}]ac.north west) {Added by randomization};
\draw[d2edgeb] (2.28,-2.05) -- node[right,font=\tiny\sffamily] {4 of 12} (2.28,-3.20);
\draw[d2edged] (7.33,-2.05) -- node[right,font=\tiny\sffamily] {4 of 12} (7.33,-3.20);
\draw[d2edgeb] (12.35,-2.05) -- node[right,font=\tiny\sffamily] {new} (12.35,-3.20);
\node[d2mid,text width=60mm,anchor=north west] at (4.20,-7.35)
  {Phase 1 v1.0.0 June 21 and Phase 2 v1.1.0 June 23, two days apart};
\pnote{0}{-8.55}{The three lower containers are the same width and start at the same y, so the
sort reads as a partition of one\\set. The inter-container gutter is 0.50 cm and
the inter-leaf gutter 0.20 cm, which is what makes the nesting legible.}
\end{appfig}
\vspace{-0.6cm}
\figcaption{Figure 12. The twelve Phase 1 protocol sections sorted into what
Phase 2\\carried unchanged, what it replaced, and what randomization made
necessary.}
\end{appfloat}

\begin{apptable}
\begin{tabularx}{\textwidth}{>{\raggedright\arraybackslash}p{2.85cm} >{\raggedright\arraybackslash}p{1.35cm} >{\raggedright\arraybackslash}p{1.85cm} >{\raggedright\arraybackslash}p{1.55cm} >{\raggedright\arraybackslash}p{1.30cm} >{\raggedright\arraybackslash}X}
\headrow \thead{Document} & \thead{Version} & \thead{Deposit} & \thead{Characters} & \thead{Figures} & \thead{Elapsed against the prior system}\\
\midrule
Phase 1 protocol & v1.0.0 & Jun 21, 2026 & 169,878 & 20 & 2 days against 6 to 12 months \\
Phase 2 protocol & v1.1.0 & Jun 23, 2026 & 199,421 & 22 & 2 days against 6 to 12 months \\
Both, combined & Two & Jun 21 to 23 & 369,299 & 42 & 4 days against a full design cycle \\
\end{tabularx}
\end{apptable}
\vspace{-0.6cm}
\tabcap{Table 13. Both protocol deposits, their versions, character counts and\\figure
counts, and the elapsed production time against the prior system.}

\subsection{The site the protocol assumes}

Figure 13 draws the equipment a participant is operated within, and the one
boundary that makes the protocol's on-premises claim checkable. Four layers sit
inside the building: oversight, control, device, and record. One node sits
outside it. The prohibited path from the on-premises model host to any external
network is the only struck-through edge in this paper, and a site that cannot
demonstrate that edge struck is not running this protocol, whatever else it
satisfies.

% ---------------------------------------------------------------------------
% FIGURE 13.  diagrams (python)-type, clustered infrastructure.  Four stacked
% clusters at a uniform 2.65 cm pitch, one long-dashed building boundary, one
% struck-through prohibited path.
% ---------------------------------------------------------------------------
\begin{appfloat}[!tb]
\begin{appfig}[diagrams (python)-type / clustered infrastructure]
\dgnode{dgtile}{1.55}{0}{\glyphuser}{Attending surgeon}{o1}
\dgnode{dgtile}{5.15}{0}{\glyphteam}{Study coordinator}{o2}
\dgnode{dgtile}{8.75}{0}{\glyphmon}{Independent monitor}{o3}
\begin{scope}[on background layer]\node[dgcluster2,fit=(o1)(o1l)(o2)(o2l)(o3)(o3l)] (l1) {};\end{scope}
\node[dgctitle2,anchor=south west] at ([yshift=1.5mm]l1.north west) {Layer 1, oversight};
\dgnodew{dgtiled}{1.55}{-2.65}{\glyphai}{On-premises LLM host}{k1}
\dgnodew{dgtiled}{4.35}{-2.65}{\glyphgear}{Advisory arbiter}{k2}
\dgnodew{dgtiled}{7.15}{-2.65}{\glyphsignal}{Heartbeat bus at 10 kHz}{k3}
\dgnodew{dgtiled}{9.95}{-2.65}{\glyphstop}{Emergency stop under 3 ms}{k4}
\begin{scope}[on background layer]\node[dgcluster,fit=(k1)(k1l)(k2)(k2l)(k3)(k3l)(k4)(k4l)] (l2) {};\end{scope}
\node[dgctitle,anchor=south west] at ([yshift=1.5mm]l2.north west) {Layer 2, control};
\dgnode{dgtilem}{1.55}{-5.30}{\glyphrobot}{Eight-arm robotic platform}{v1}
\dgnode{dgtilem}{4.35}{-5.30}{\glyphhand}{Per-arm force cap 3 N}{v2}
\dgnode{dgtilem}{7.15}{-5.30}{\glyphhand}{Cross-arm cap 18 N}{v3}
\dgnode{dgtilem}{9.95}{-5.30}{\glyphshield}{Vascular no-fly gating}{v4}
\begin{scope}[on background layer]\node[dgcluster,fit=(v1)(v1l)(v2)(v2l)(v3)(v3l)(v4)(v4l)] (l3) {};\end{scope}
\node[dgctitle,anchor=south west] at ([yshift=1.5mm]l3.north west) {Layer 3, device};
\dgnodeg{dgtileg}{1.55}{-7.95}{\glyphdb}{Operative log store}{c1}
\dgnodeg{dgtileg}{5.15}{-7.95}{\glyphchart}{Pathology and imaging}{c2}
\dgnodeg{dgtileg}{8.75}{-7.95}{\glyphlock}{Verification hash registry}{c3}
\begin{scope}[on background layer]\node[dgcluster2,fit=(c1)(c1l)(c2)(c2l)(c3)(c3l)] (l4) {};\end{scope}
\node[dgctitle2,anchor=south west] at ([yshift=1.5mm]l4.north west) {Layer 4, record};
% boundary
\draw[trialburg,line width=1.2pt,dash pattern=on 6pt off 3pt] (11.85,1.05) -- (11.85,-9.05);
\node[dgctitle,rotate=90,anchor=south] at (11.75,-4.00) {site boundary};
\dgnodeg{dgtileg}{13.55}{-2.65}{\glyphcloud}{Any external network}{ext}
% chain edges
\draw[dgedgeb] (o1l.south) -- node[right,font=\tiny\sffamily] {oversees} (k1.north);
\draw[dgedgeb] (k1l.south) -- node[right,font=\tiny\sffamily] {advises} (v1.north);
\draw[dgedgeb] (v1l.south) -- node[right,font=\tiny\sffamily] {writes} (c1.north);
\draw[dgedged] (c3.north) -- node[right,font=\tiny\sffamily] {attests} (o3l.south);
% prohibited path
\draw[trialchar,line width=1pt,dotted] (k1.east) -- (11.85,-2.65);
\pxmark{11.85}{-2.65}
\node[font=\tiny\sffamily,text=trialburg,anchor=south] at (10.35,-2.45) {no path};
\legkey{0.35}{-9.55}{trialburgp}{inside the building}
\legkey{4.20}{-9.55}{trialmistl}{record and external}
\pnote{0.35}{-10.25}{The four clusters sit at a uniform 2.65 cm pitch, 1.05 cm more than the tallest
tile and label pair, so no frame\\touches another. The attestation edge rises
7.95 cm through the corridor at x = 8.75, which carries no tile.}
\end{appfig}
\vspace{-0.6cm}
\figcaption{Figure 13. The four-layer on-premises site stack, and the single
boundary\\across which no inference request, model weight, or patient record
passes.}
\end{appfloat}

The boundary is not a preference. It is what makes the consent language
operable: a participant told that an on-premises model advises the operation can
verify the claim by asking what crosses the building line, and a monitor can
audit it by inspecting one network configuration rather than a vendor's
assurances \cite{onpremwhippl, clintrust}. The same boundary is what the
autonomy disclosure duty in \S5 attaches to.

\subsection{Why two protocols rather than one}

A Phase 1 protocol and a Phase 2 protocol two days apart is an unusual pairing,
and the reason is structural rather than opportunistic. A first-in-human study
of a device that advises an operation has to answer a question the drug arm
alone cannot answer: whether the advisory changes the operation in a way that a
randomized comparison would later detect. Writing the randomized study first, as
a target, forces the first-in-human study to collect what the randomized study
will need. The Phase 2 document therefore fixes the primary endpoint,
progression-free survival per RECIST 1.1 with blinded independent central review
as a sensitivity analysis, and the Phase 1 document collects the operative and
pathology inputs that endpoint depends on \cite{onpremwhippl, phase2proto}.

The pairing also exposes what the prior system's sequencing hides. In the prior
system the Phase 2 protocol is written after the Phase 1 result is known, which
means the Phase 1 case report form was designed without knowing which of its
fields would matter. Here both documents were written against one another, and
Figure 12 shows the consequence: only four of the twelve Phase 1 sections had to
be replaced when the study randomized, and none of the four is a data-collection
section. Compliance, introduction, discontinuation and additional considerations
carried unchanged. Summary, objectives, design and statistics were replaced,
which is exactly the set a change from single-arm dose-finding to randomized
comparison should replace and no more.

\subsection{What the two protocols assume about the site}

Both protocols assume a site that does not yet exist in the form they require,
and both say so. The Phase 1 protocol assumes one qualified high-volume
hepato-pancreato-biliary center with an eight-arm robotic platform, an
on-premises model host inside the building boundary of Figure 13, a pathology
service able to adjudicate margin status at day 90, and an institutional review
board willing to review a combined IND and IDE submission with a Physical AI
consent opt-out \cite{onpremwhippl}. The Phase 2 protocol assumes eight such
centers, plus central masked pathology, blinded endpoint adjudication, and a
central randomization service with four stratification factors
\cite{phase2proto}.

None of those instruments can be obtained by a sponsor of 2.6 full-time
equivalents acting alone, which is the point \S6 and \S8 return to. The
protocols are written to be executable by an institution that has them, not to
be executable by their author. That is a deliberate limit on what document
production can accomplish, and stating it here is more useful than implying
otherwise: the documents remove the drafting bottleneck and they do not remove
the institutional one.

\subsection{The safety architecture in one paragraph}

A reader who takes only one thing from \S\ref{sec:protocol} should take the
safety architecture, because it is the part a clinician will test first. Four
independent limits act on the operation at once. A heartbeat broadcast at 10 kHz
means the control layer knows within 100 microseconds that an arm has stopped
reporting. A cross-arm emergency stop budget at or below 3 ms means that when
any arm asserts a stop, every other arm has halted before a human could have
noticed the condition. Per-arm tip-force caps at or below 3 N and a cumulative
cross-arm cap at or below 18 N mean the platform cannot apply more force than
the caps allow even if every controller above them is wrong. Vascular no-fly
gating around the superior mesenteric and portal veins means the geometric
region in which those forces would be most dangerous is excluded before any of
the other three limits is tested \cite{onpremwhippl}.

Those four are cyber-physical rather than advisory: none depends on the model
being correct, and none depends on the surgeon reacting in time. The model's
role sits above them, in the perioperative restart advisory and the operative
guidance, and Figure 10's guard on the transition out of the operative state is
written against the force caps rather than against the model's confidence. That
ordering is the protocol's answer to the trust question \S\ref{sec:legislation}
raises: the clinician is asked to rely on the model for advice and on the
hardware for safety, and the two reliances are separable by design.

\subsection{What the two protocols would need to open}

The instruments the protocols require are the same class as those the IND
requires, and the register in \S\ref{sec:limitations} lists all seven together.
For the Phase 1 study the binding ones are a qualified high-volume
hepato-pancreato-biliary center, an operating room with the eight-arm platform
installed and accepted, a pathology service able to adjudicate margin status at
day 90, and an institutional review board of record. For the Phase 2 study the
binding ones are eight such centers plus central masked pathology, blinded
endpoint adjudication, and a central randomization service.

The protocols are deliberately not written down to what one sponsor could run
alone. A single-site, unblinded, unadjudicated study would have been easier to
open and would not have answered the comparative question, and the Phase 1
document says as much when it defers that question to a randomized trial
\cite{onpremwhippl, phase2proto}. Writing to the study that answers the question
and then stating plainly what is missing is more useful to a funder than writing
to the study that could be opened tomorrow.
